% ==============================================================================
% estilo.tex — Preámbulo LaTeX común de mis PDF
%
% QUÉ ES
%   Código LaTeX que se agrega al "preámbulo" (la parte de arriba) de cada
%   documento, antes de que empiece el texto. Define el aspecto: fuentes,
%   emojis, listas, tablas, cuadros y bloques de código.
%
% CÓMO SE USA
%   No se llama a mano: lo agrega al preámbulo el filtro
%   filters/preambulo.lua (que a su vez está activado en estilo.yaml).
%   Se usa con:
%       pandoc -d estilo nota.md -o nota.pdf
%
% CÓMO SE EDITA
%   Cada sección (0 a 9) es independiente. Para sacar algo del estilo, borrá
%   o comentá la sección entera. En LaTeX, un comentario empieza con "%".
%   Para ver el LaTeX completo que genera pandoc (con este archivo ya
%   insertado en el preámbulo), sin compilar el PDF:
%       pandoc -d estilo nota.md -s -t latex -o nota.tex
%
% REQUISITOS (Debian/Ubuntu)
%   sudo apt install pandoc texlive-luatex texlive-latex-recommended \
%        texlive-latex-extra texlive-fonts-recommended lmodern \
%        fonts-noto-color-emoji fonts-dejavu-core fonts-dejavu-mono
%   - texlive-luatex:         el motor lualatex (secciones 0 y 5 lo necesitan).
%   - texlive-latex-extra:    trae fvextra (sección 4).
%   - fonts-noto-color-emoji: los emojis a color (sección 0).
%   - fonts-dejavu-core:      respaldo de símbolos en texto (sección 0).
%   - fonts-dejavu-mono:      respaldo de símbolos y árboles en código (sección 0).
%
% SECCIONES
%   0) Fuente principal y emojis (se cambia con "mainfont" en la nota)
%   1) Listas con más aire
%   2) Tablas con más aire
%   3) Cuadros para las citas (> en Markdown)
%   4) Código: cortar las líneas largas
%   5) Palabras largas sin espacios (URLs, hashes)
%   6) \needspace (evitar títulos solos al pie de la página)
%   7) Ancho máximo de las imágenes (se cambia con "ancho-imagenes")
%   8) Figuras con epígrafe fijas donde están en el Markdown
%   9) Epígrafes sin "Figura N:"
% ==============================================================================


% ------------------------------------------------------------------------------
% 0) Fuentes principales embebidas con respaldo: emojis a color y símbolos
%
%    Las tres fuentes estándar del repositorio están incluidas directamente en
%    toolchain/pandoc/fonts/ para no depender del gestor de paquetes de cada
%    distribución Linux:
%      - Cuerpo (mainfont):  Libertinus Serif
%      - Títulos (sansfont): Inter
%      - Código (monofont):  JetBrains Mono
%
%    Además, con lualatex se configura una cascada de respaldo cuando falta un
%    carácter:
%      1.º Noto Color Emoji  (emojis a color; "mode=harf")
%      2.º DejaVu Sans       (símbolos matemáticos/técnicos como ✓ o →)
%    Los bloques de código tienen su propio respaldo (primero DejaVu Sans Mono
%    para árboles ├─ │ └ y flechas ↕ en blanco y negro, y segundo Noto Color
%    Emoji por si hay un emoji puro dentro de `código`).
% ------------------------------------------------------------------------------
\ifdefined\fuenteprincipal\else
  \def\fuenteprincipal{Libertinus Serif}
\fi
\ifdefined\fuentesans\else
  \def\fuentesans{Inter}
\fi
\ifdefined\fuentemono\else
  \def\fuentemono{JetBrains Mono}
\fi
\def\nombreLibertinus{Libertinus Serif}
\def\nombreInter{Inter}
\def\nombreJetBrains{JetBrains Mono}
\ifLuaTeX
  \directlua{
    luaotfload.add_fallback("emojis",
      {"Noto Color Emoji:mode=harf;", "DejaVu Sans:"})
    luaotfload.add_fallback("monoletras",
      {"DejaVu Sans Mono:", "Noto Color Emoji:mode=harf;"})
  }
  \ifdefined\carpetaFuentes
    \ifx\fuenteprincipal\nombreLibertinus
      \setmainfont{libertinus-serif}[
        Path           = {\carpetaFuentes/libertinus-serif/},
        Extension      = .ttf,
        UprightFont    = *,
        ItalicFont     = *-italic,
        BoldFont       = *-700,
        BoldItalicFont = *-700-italic,
        RawFeature     = {fallback=emojis}
      ]
    \else
      \setmainfont{\fuenteprincipal}[RawFeature={fallback=emojis}]
    \fi
    \ifx\fuentesans\nombreInter
      \setsansfont{inter}[
        Path           = {\carpetaFuentes/inter/},
        Extension      = .ttf,
        UprightFont    = *,
        ItalicFont     = *-italic,
        BoldFont       = *-700,
        BoldItalicFont = *-700-italic,
        RawFeature     = {fallback=emojis}
      ]
    \else
      \setsansfont{\fuentesans}[RawFeature={fallback=emojis}]
    \fi
    % Letra del código con respaldo: JetBrains Mono / Latin Modern Mono (y
    % muchas otras) no traen todos los trazos de árboles (├ └ │ ─) o flechas
    % como ↕; esos salen de DejaVu Sans Mono. Si la nota pone "monofont: ...",
    % se usa esa letra.
    \ifx\fuentemono\nombreJetBrains
      \setmonofont{jetbrains-mono}[
        Path           = {\carpetaFuentes/jetbrains-mono/},
        Extension      = .ttf,
        UprightFont    = *,
        ItalicFont     = *-italic,
        BoldFont       = *-700,
        BoldItalicFont = *-700-italic,
        RawFeature     = {fallback=monoletras}
      ]
    \else
      \setmonofont{\fuentemono}[RawFeature={fallback=monoletras}]
    \fi
  \else
    \setmainfont{\fuenteprincipal}[RawFeature={fallback=emojis}]
    \setsansfont{\fuentesans}[RawFeature={fallback=emojis}]
    \setmonofont{\fuentemono}[RawFeature={fallback=monoletras}]
  \fi
\else
  \ifdefined\setmainfont
    \setmainfont{\fuenteprincipal}
    \setsansfont{\fuentesans}
    \setmonofont{\fuentemono}
  \fi
\fi


% ------------------------------------------------------------------------------
% 1) Listas con más aire entre ítems
%
%    La separación se define con \setlist (del paquete enumitem) y vale
%    para todas las listas, con o sin líneas en blanco entre los ítems en el
%    Markdown: así una lista sale igual en el PDF de las dos formas.
%      - Ítems:     0,55 em entre uno y otro (aprox. media línea).
%      - Subítems (nivel 2 en adelante): 0,8 em entre uno y otro, y 0,8 em
%        de aire antes y después de la sublista (topsep), para que no quede
%        pegada al texto del ítem de arriba ni al ítem de abajo. OJO: LaTeX
%        usa el mismo valor arriba y abajo; no se puede cambiar uno solo.
%      - Antes y después de la lista: 0,18 em (topsep del nivel 1). Parece
%        poco, pero se suma al espacio que ya deja el párrafo: así, del
%        párrafo al primer ítem hay la misma distancia que del ítem a su
%        subítem, y lo mismo del último ítem (o subítem) al párrafo de abajo.
%        Si cambiás el 0.8em de los subítems, ajustá también este número.
%    parsep=0pt evita que se sume además el espacio entre párrafos.
%
%    Pandoc pone \tightlist en las listas "apretadas" (sin líneas en blanco
%    entre ítems) para dejarlas con separación 0. Se anula (queda vacío),
%    porque si no pisaría los valores de acá.
%
%    Para más o menos aire, cambiar los números de los dos \setlist.
% ------------------------------------------------------------------------------
\usepackage{enumitem}
\setlist{itemsep=0.55em, parsep=0pt}
\setlist[1]{topsep=0.18em, partopsep=0pt}
\setlist[2,3,4]{itemsep=0.8em, topsep=0.8em, partopsep=0pt}
\renewcommand{\tightlist}{}


% ------------------------------------------------------------------------------
% 2) Tablas con más aire entre filas
%
%    \arraystretch multiplica la altura de las filas (el valor por defecto es
%    1; acá 1,5). Se aplica solo dentro de las tablas ("longtable", que es lo
%    que genera pandoc). Si se pusiera global, también ensancharía las líneas
%    de autor y fecha del título.
% ------------------------------------------------------------------------------
\AddToHook{env/longtable/begin}{\renewcommand{\arraystretch}{1.5}}


% ------------------------------------------------------------------------------
% 3) Cuadros para las citas: "Idea clave", "Ojo", "Para el examen"...
%
%    Todo bloque de cita de Markdown (las líneas que empiezan con ">") se
%    dibuja como un cuadro con una barra azul a la izquierda y fondo celeste
%    pastel, con aire arriba y abajo. OJO: esto cambia TODAS las citas, no se
%    puede tener a la vez una cita "normal" y un cuadro.
%
%    Para cambiar los colores, editar los dos \definecolor (formato RGB, de 0
%    a 255). El resto arma el cuadro en tres pasos: mide el contenido en una
%    caja (\cuadrobox), dibuja la barra con la altura de esa caja, y pone la
%    caja sobre el fondo con 10 pt de margen a cada lado.
%    Solo usa LaTeX base y color (no necesita paquetes extra).
% ------------------------------------------------------------------------------
\definecolor{barcolor}{RGB}{74,111,165}
\definecolor{fondocolor}{RGB}{232,241,251}
\newsavebox{\cuadrobox}
\renewenvironment{quote}{%
  \par\addvspace{10pt}%
  \begin{lrbox}{\cuadrobox}%
  \begin{minipage}[t]{\dimexpr\linewidth-24pt\relax}%
  \setlength{\parskip}{0.5em}}%
  {\par\end{minipage}\end{lrbox}%
  \par\noindent\hbox{%
    {\color{barcolor}\vrule width 2.5pt
      height \dimexpr\ht\cuadrobox+6pt\relax
      depth \dimexpr\dp\cuadrobox+6pt\relax}%
    {\setlength{\fboxsep}{0pt}%
     \colorbox{fondocolor}{\vrule width 0pt
      height \dimexpr\ht\cuadrobox+6pt\relax
      depth \dimexpr\dp\cuadrobox+6pt\relax
      \hskip 10pt\usebox{\cuadrobox}\hskip 10pt}}}%
  \par\addvspace{10pt}}


% ------------------------------------------------------------------------------
% 4) Bloques de código: las líneas largas se cortan
%
%    Sin esto, una línea de código más ancha que la hoja se sale por la
%    derecha. Con el paquete fvextra, se corta y la línea que continúa empieza
%    con una flechita (↪).
%      breaklines            corta las líneas largas
%      breakanywhere         permite cortar también dentro de una palabra larga
%      breaknonspaceingroup  hace que ese corte funcione también dentro de los
%                            textos con color del resaltado de sintaxis (por
%                            ejemplo, un string muy largo)
%    Las dos redefiniciones de abajo hacen que lo anterior valga para ambos
%    tipos de bloque de pandoc: "verbatim" (los que no tienen lenguaje) y
%    "Highlighting" (los que sí tienen, con colores). Este último solo existe
%    si la nota tiene código con lenguaje, por eso se pregunta antes con
%    \ifdefined.
% ------------------------------------------------------------------------------
\usepackage{fvextra}
\fvset{breaklines, breakanywhere, breaknonspaceingroup}
\RecustomVerbatimEnvironment{verbatim}{Verbatim}{}
\ifdefined\Highlighting
  \RecustomVerbatimEnvironment{Highlighting}{Verbatim}{commandchars=\\\{\}}
\fi


% ------------------------------------------------------------------------------
% 5) Palabras largas sin espacios (URLs, hashes, rutas) en tablas, listas
%    o párrafos
%
%    Una palabra muy larga y sin espacios no tiene dónde cortarse y se sale de
%    la hoja (típico en las celdas de una tabla). Este bloque le avisa a TeX
%    que, SOLO en las palabras de 20 caracteres o más seguidos, puede cortar
%    entre dos letras. Es el último recurso: se pone un "corte muy caro"
%    (penalización 9000 de 10000) para que TeX lo use únicamente si no hay
%    otra forma de acomodar la línea. Las palabras normales no se tocan, y
%    tampoco las fórmulas matemáticas. Solo funciona con lualatex.
%
%    Para cambiar desde cuántos caracteres se corta, editar el número de
%    "local MINIMO". Con menos de 20, empiezan a cortarse palabras comunes.
%
%    CUIDADO al editar el código de Lua de abajo: TeX lo lee primero, así que
%      - no se pueden usar comentarios de Lua (--): TeX junta todas las líneas
%        en una sola y el comentario taparía el resto del código;
%      - no se pueden usar los símbolos de porcentaje, numeral ni virgulilla
%        (tampoco el operador "no igual" de Lua): TeX los interpretaría antes.
%    Por eso las explicaciones van acá afuera.
% ------------------------------------------------------------------------------
\ifLuaTeX
  \directlua{
    local GLYPH = node.id("glyph")
    local KERN = node.id("kern")
    local PENALTY = node.id("penalty")
    local MATH = node.id("math")
    local MINIMO = 20
    local function cortar(head)
      local enmat = false
      local n = head
      while n do
        if n.id == MATH then
          enmat = (n.subtype == 0)
          n = n.next
        elseif n.id == GLYPH and not enmat then
          local largo = 0
          local m = n
          while m and (m.id == GLYPH or (m.id == KERN and m.subtype == 0)) do
            if m.id == GLYPH then largo = largo + 1 end
            m = m.next
          end
          if largo >= MINIMO then
            local c = n
            while c and not (c == m) do
              if c.id == GLYPH then
                local sig = c.next
                if sig and sig.id == KERN then sig = sig.next end
                if sig and not (sig == m) and sig.id == GLYPH then
                  local p = node.new(PENALTY)
                  p.penalty = 9000
                  head, c = node.insert_after(head, c, p)
                end
              end
              c = c.next
            end
          end
          n = m
        else
          n = n.next
        end
      end
      return head
    end
    luatexbase.add_to_callback("pre_linebreak_filter", cortar, "cortar-palabras-largas")
  }
\fi


% ------------------------------------------------------------------------------
% 6) \needspace{altura}: evitar títulos solos al pie de la página
%
%    Uso: en el .md, en una línea suelta justo antes de un título:
%        \needspace{5cm}
%    Si quedan menos de 5 cm libres en la página, salta a la siguiente. Sirve
%    para que un título no quede abandonado al final con su texto en la página
%    de abajo. Es la definición del paquete "needspace", copiada acá para no
%    tener que cargar otro paquete. \providecommand la define solo si no existe.
% ------------------------------------------------------------------------------
\makeatletter
\providecommand{\needspace}[1]{\begingroup\setlength{\dimen@}{#1}%
  \vskip\z@\@plus\dimen@\penalty-100\vskip\z@\@plus-\dimen@
  \vskip\dimen@\penalty9999\vskip-\dimen@\vskip\z@skip\endgroup}
\makeatother


% ------------------------------------------------------------------------------
% 7) Ancho máximo de las imágenes
%
%    Pandoc ya achica las imágenes que no entran en el ancho del texto. Esto
%    baja ese tope: ninguna imagen pasa de \anchoimagenes (por ejemplo, el
%    80 % del ancho del texto). Las imágenes más chicas quedan como están, y
%    se mantiene la proporción.
%
%    El valor se elige con "ancho-imagenes" en metadatos.yaml (para todas las
%    notas) o en la cabecera de una nota (solo esa). El filtro
%    filters/preambulo.lua lo convierte en \anchoimagenes. Sin ese valor, el
%    tope es el ancho completo del texto, como hace pandoc.
%
%    Una imagen con su propio tamaño en el Markdown, por ejemplo
%        ![](foto.png){width=40%}
%    usa ese tamaño, sin importar este tope.
%
%    Pandoc cambió la forma de achicar imágenes en la versión 3.2.1:
%      - hasta la 3.2: usa \maxwidth (Debian 13 trae la 3.1.11);
%      - desde la 3.2.1: usa \pandocbounded.
%    Acá se redefinen los dos, copiando lo que hace pandoc pero con
%    \anchoimagenes en lugar del ancho completo (\linewidth), así funciona con
%    cualquier versión. Si pandoc vuelve a cambiar el nombre, este tope deja
%    de aplicarse sin dar error: las imágenes vuelven al ancho completo.
% ------------------------------------------------------------------------------
\providecommand{\anchoimagenes}{\linewidth}
\makeatletter
\def\maxwidth{\ifdim\Gin@nat@width>\anchoimagenes\anchoimagenes\else\Gin@nat@width\fi}
\ifdefined\pandocbounded
  \renewcommand*\pandocbounded[1]{%
    \sbox\pandoc@box{#1}%
    \Gscale@div\@tempa{\textheight}{\dimexpr\ht\pandoc@box+\dp\pandoc@box\relax}%
    \Gscale@div\@tempb{\anchoimagenes}{\wd\pandoc@box}%
    \ifdim\@tempb\p@<\@tempa\p@\let\@tempa\@tempb\fi
    \ifdim\@tempa\p@<\p@\scalebox{\@tempa}{\usebox\pandoc@box}%
    \else\usebox{\pandoc@box}%
    \fi}
\fi
\makeatother


% ------------------------------------------------------------------------------
% 8) Figuras con epígrafe fijas donde están en el Markdown
%
%    Una imagen con epígrafe, ![Texto](foto.png), es una "figura", y LaTeX la
%    trata como flotante: la pone donde le queda mejor (por ejemplo, arriba de
%    la página siguiente) y sube el texto de abajo a ocupar el hueco. Así la
%    figura puede terminar lejos del párrafo que la menciona.
%
%    El paquete float agrega la posición "H" (exactamente acá), y
%    \floatplacement la pone por defecto en todas las figuras. Si una figura
%    no entra en lo que queda de la página, pasa entera a la siguiente y deja
%    un espacio en blanco al final de la anterior.
%
%    Las imágenes sin epígrafe no son figuras: siempre quedan en su lugar,
%    con o sin esto. Las tablas tampoco se ven afectadas (pandoc no las hace
%    flotantes).
%
%    Para que una nota vuelva al comportamiento normal de LaTeX, en su
%    cabecera:
%        header-includes: \floatplacement{figure}{tbp}
%    El paquete float viene en texlive-latex-recommended.
% ------------------------------------------------------------------------------
\usepackage{float}
\floatplacement{figure}{H}


% ------------------------------------------------------------------------------
% 9) Epígrafes sin "Figura N:"
%
%    Una figura, ![Consecuencias primarias](foto.png), sale normalmente con
%    el epígrafe "Figura 1: Consecuencias primarias", y el número sigue en
%    cada figura. Con esto sale solo el texto: "Consecuencias primarias".
%    Solo cambia las figuras; las tablas con título no se tocan.
%
%    Para que una nota vuelva a numerar las figuras, en su cabecera:
%        header-includes: \captionsetup[figure]{labelformat=simple}
%    El paquete caption viene en texlive-latex-recommended (pandoc ya lo
%    carga en muchas notas; cargarlo otra vez no hace nada).
% ------------------------------------------------------------------------------
\usepackage{caption}
\captionsetup[figure]{labelformat=empty}
